\chapter{Constants: provenance and machine verification}\label{sec-appendix}
% Ported verbatim from paper/top3_joint_moment_bound.tex lines 1417-1454.
% Only cross-reference wording was adapted; no mathematics was altered.

A reader cannot tell from the statements alone which of the many integer
coefficients in this paper are forced by an interpolation problem and which
were chosen and then verified. We record the distinction.

\textbf{Uniquely determined.}
(i)~The polynomial $f$ of \S6 is the unique polynomial of degree at most
four with double contact to $|w|+w$ at $w=-3/8$ and at $w=1/2$ and a contact
at $w=1$; five conditions determine five coefficients, so once the nodes are
fixed the integers in $f$ (including $41503=7^3\cdot11^2$) are forced.
(ii)~$z_*$ of \S8 is determined by the same nodes together with the
requirement that the cubic term vanish, i.e.\ by the factor $(t+9/8)$.
(iii)~$\cstar$ of Lemma~\ref{attainment} is the value of an explicit
one-variable function at the critical point $b=(7+5\sqrt7)/126$.
(iv)~The four rational certificates of \S7 and \S9 are obtained by
substituting the stated rational bounds for $\sqrt5$ and $\sqrt3$ into
$d_0,d_1,d$; no search is involved.

\textbf{Chosen, then verified exactly.}
(i)~The nodes $-3/8,\,1/2,\,1$ themselves, and the majorant $h_0$ taken
from~\cite[Lemma 4.4]{STW}.
(ii)~The parameters $\ell=937/1400$, $\eta=1/10$, $\theta=6/25$ of \S7 and
the square coefficients of \eqref{squares} and \eqref{Q8}; these were
selected so that all residual coefficients remain positive, and every
positivity claim is then checked by exact rational arithmetic.
(iii)~The divisor $125$ and $\kappa=3/125000$ of \eqref{Cstar8}, and the
constant $3/2$ of Lemma~\ref{contact-square}; the quantified loss is
discussed in Remark~\ref{rem-theta}.

\textbf{Verification protocol.} Every displayed identity, factorization,
completion of squares and ``$>0$'' assertion was re-checked in exact
arithmetic: rationals as exact fractions and $\sqrt3,\sqrt5$ in the
quadratic fields $\mathbb Q(\sqrt3)$, $\mathbb Q(\sqrt5)$, never in floating
point. Pointwise inequalities were additionally scanned on a
$2\times10^6+1$ point grid, and the finite-graph transfer was checked on all
$33864$ graphs of order $3\le n\le6$. The verification scripts are listed in
the accompanying audit report of 2 November 2026.


\section*{Additions made in this thesis}
\addcontentsline{toc}{section}{Additions made in this thesis}
Two constants and one exact spectral identity appear in
Chapter~\ref{ch-numerical} and Section~\ref{sec-weighted} and do not occur in
the preprint~\cite{Ji}; their provenance is recorded here for completeness.

\textbf{Uniquely determined.}
(v)~The equal-blow-up coefficient $(\sqrt5+\sqrt6+3)/6$ of
Proposition~\ref{equal-blowup} is forced: it is $\Lambda_3(K_1\vee C_5)+3$
divided by $6$, and both $\Lambda_3(K_1\vee C_5)=\sqrt5+\sqrt6$ and the shift
$+3$ from the clique parts follow from Lemma~\ref{blowup-spec} without any
choice.
(vi)~The optimal rim weight $b^*=(7+5\sqrt7)/126$ is the critical point of
the one-variable quotient function computed in Lemma~\ref{attainment}, and the
apex weight $p^*=(91-25\sqrt7)/126$ equals $1-5b^*$; neither is chosen.

\textbf{Chosen, then verified.}
(iv)~The orders $n\le7$ for the exhaustive search, the blow-up sizes
$q\in\{1,2,3,5,10,20,50,100,250,500\}$, and the weighted sample sizes
$4000$ per order in $n\in\{4,6,9,15,30,60\}$. These are sampling decisions;
they affect the strength of the numerical evidence and nothing else.

\textbf{Verification protocol for the additions.} The exhaustive searches
enumerated all $2^{n(n-1)/2}$ labelled graphs, computed
$\lambda_1+\lambda_2+\lambda_3$ with a symmetric eigensolver in double
precision, and compared against $cn-3$ evaluated from the exact expression
for $c$; the margins observed are at least $2.2\times10^{-2}$, which is seven
orders of magnitude above the double-precision error of a $7\times7$ symmetric
eigenvalue computation. The number of labelled maximisers found at $n=6$ and
$n=7$ is $72=6!/10$ and $252=7!/20$, matching
$|\operatorname{Aut}(K_1\vee C_5)|=10$ and
$|\operatorname{Aut}(K_2\vee C_5)|=2\cdot10=20$, which is the consistency
check for uniqueness up to isomorphism. The exact identities
$(\sqrt5+\sqrt6+3)/6\cdot6q-3=\Lambda_3(H[q])$ were re-verified in exact
arithmetic in $\mathbb Q(\sqrt5,\sqrt6)$ for $q\le10$ and in double precision
for all $q$ in the table.
